\documentclass[11pt,a4paper]{article}
\usepackage[margin=2.1cm]{geometry}
\usepackage{amsmath,amssymb}
\usepackage{booktabs}
\usepackage{enumitem}
\usepackage{hyperref}
\usepackage{xcolor}
\usepackage{tcolorbox}
\usepackage{fancyhdr}

\pagestyle{fancy}
\fancyhead[L]{\small CS5491 AI --- Course Project}
\fancyhead[R]{\small Resource-Aware Configuration Search}
\setlength{\parindent}{0pt}
\setlength{\parskip}{0.22em}
\setlist[itemize]{itemsep=0pt,topsep=1pt}

\title{\vspace{-2.2cm}\textbf{Resource-Aware Heuristic Search for\\Edge Vision Deployment Configuration Selection}\\[4pt]
\large Final Report (draft)}
\author{WU JUNYU \quad 72655994 \quad 72655994@cityu-dg.edu.cn}
\date{CS5491 Artificial Intelligence}

\begin{document}
\maketitle
\vspace{-1.2em}

\begin{tcolorbox}[colback=orange!6,colframe=orange!55!black,title=\textbf{Status of the numbers in this draft}]
Every numerical table below was produced by the \textbf{pipeline verification run}
on \texttt{benchmarks/synthetic\_demo.csv}, whose rows carry
\texttt{source=synthetic}. They demonstrate that the code produces the required
artefacts; they are \emph{not} measurements and must be replaced by the
Raspberry Pi table before submission. The commands and the section structure are
final.
\end{tcolorbox}

\section{Introduction}

Deploying a vision model on an edge device is a three-way trade-off between
detection accuracy, inference latency and memory. Which point of that trade-off
you obtain is determined by the \emph{deployment configuration}: the detector
variant, the input resolution, the numerical precision and the number of CPU
threads. Exhaustively benchmarking every candidate is acceptable for a small
space, but the cost grows quickly, and on a device like a Raspberry Pi each
measurement is slow.

We therefore treat configuration selection as a \textbf{resource-constrained
search problem}. Each configuration is an expensive black-box query returning an
accuracy, a latency and a memory value; only a fixed number of queries is
affordable; the goal is a feasible configuration of high accuracy.

Our research question is:

\begin{quote}\itshape
Given a limited number of configuration evaluations, does an analytic,
resource-aware score select better deployment configurations than random or
purely data-driven local search?
\end{quote}

Following the experimental protocol of the course topic \emph{Neural Spectral
Capacity}~\cite{nsc}, the score is \emph{analytic}: it is computed from the
configuration specification alone, without running the detector. The topic's
protocol --- cheap score, rank correlation, fixed-budget selection, ablation ---
transfers directly; only the objects being scored change, from Transformer
architectures to deployment configurations.

\section{Method}

\subsection{Configuration space}

A configuration is a point in
$\{\text{YOLOv5n},\text{YOLOv5s},\text{YOLOv5m}\} \times \{320,416,640\}
\times \{\text{FP32},\text{INT8}\} \times \{1,2,4\}$, giving 54 candidates
(27 when restricted to FP32). The space is defined in
\texttt{configs/search\_space.yaml} and enumerated deterministically, so the
search cannot drift away from the specification.

Including three model variants is deliberate. The capacity score is an
\emph{architecture} quantity: if every candidate shared one architecture, the
score would be constant and could not rank anything.

\subsection{Analytic score}

For a weight matrix $W$ with singular values $\sigma_i$, the reference capacity
is $\psi(W)=\ln\det(I+W^\top W)=\sum_i \ln(1+\sigma_i^2)$. For i.i.d.\
$\mathcal{N}(0,s^2)$ entries, the Marchenko--Pastur law gives its expectation
from the shape $(m,n)$ and the initialisation scale $s$ alone:
\begin{equation}
\mathbb{E}[\psi_{\mathrm{MP}}(m,n,s)]
= \min(m,n)\int \ln\!\big(1+s^2\max(m,n)\lambda\big)\,d\mu_{\gamma}(\lambda),
\qquad \gamma=\tfrac{\min}{\max}.
\end{equation}
The integral is evaluated numerically with the substitution
$\lambda = a + (b-a)t^2$, which removes the inverse-square-root singularity of
the MP density at the lower edge --- a plain uniform grid underestimates the
integral by up to two orders of magnitude for near-square matrices.

The architecture capacity $C(\text{model})$ sums $\psi_{\mathrm{MP}}$ over every
convolution matrix of the variant. The composite score adds two deployment
factors and the resource penalties:
\begin{equation}
\hat A(x)=C(\text{model})\left(\frac{R}{640}\right)^{\alpha} q(\text{precision}),
\qquad
S(x)=\frac{\hat A(x)}{\max \hat A}
-\lambda P_{\text{lat}}(x)-\mu P_{\text{mem}}(x),
\end{equation}
with $\alpha=1$, $q(\text{FP32})=1$, $q(\text{INT8})=0.97$, and $P_\bullet=0$
inside a constraint, growing linearly outside it.

The resource side is predicted analytically,
$\text{MACs}=\sum H_{\text{out}}W_{\text{out}}C_{\text{out}}C_{\text{in}}k^2$,
$\text{latency}=\text{MACs}/(\kappa\cdot\text{threads})+\text{overhead}$,
$\text{memory}=\text{weights}+\text{peak activation}+\text{runtime baseline}$.
The three constants are priors and are refitted from measurements by
\texttt{ResourceModel.fit}.

\subsection{Validating the architecture specification}

The analytic spec is checked against the published YOLOv5 figures
(\texttt{scripts/check\_model\_spec.py}):
\begin{center}
\begin{tabular}{@{}lcccccc@{}}
\toprule
Variant & params & ref. & ratio & GFLOPs & ref. & ratio \\
\midrule
YOLOv5n & 1.58\,M & 1.90\,M & 0.830 & 4.23 & 4.5  & 0.940 \\
YOLOv5s & 6.05\,M & 7.20\,M & 0.841 & 15.49 & 16.5 & 0.939 \\
YOLOv5m & 17.06\,M & 21.20\,M & 0.805 & 45.57 & 49.0 & 0.930 \\
\bottomrule
\end{tabular}
\end{center}
Compute matches to within 7\% with a nearly identical ratio across variants,
which is the quantity that drives the latency model. Parameters are under-counted
by about 19\% uniformly, because the detection head is modelled as a single
$1\times1$ convolution per level. That deficit is immaterial downstream: the
capacity score is scale-invariant, and the weights are a small part of the memory
estimate, which the runtime baseline dominates.

\subsection{Search algorithms}

The space is a graph in which two configurations are adjacent when they differ
in exactly one factor by one step. Four methods are compared under an identical
budget: \textbf{random} sampling; \textbf{greedy} local search that moves to the
best \emph{measured} neighbour; \textbf{heuristic} best-first search ordered by
$S(x)$; and \textbf{exhaustive} enumeration as the oracle.

The heuristic never inspects an unmeasured configuration: its ordering comes
only from $S$, and the evaluator raises an error if an unmeasured row is
requested.

\section{Experimental protocol}

All methods share one evaluator, one pool and one budget. The evaluator counts
\emph{distinct} evaluations, so a repeated query is free; this makes the reported
$n$ the quantity to compare against. Metrics are:

\begin{itemize}
  \item \textbf{Best feasible accuracy} and \textbf{regret} against the oracle;
  \item \textbf{Spearman $\rho$} and \textbf{Kendall $\tau_b$} between the score
        and the measured accuracy, with average ranks and the $\tau_b$ tie
        correction;
  \item the \textbf{accuracy--latency Pareto front}.
\end{itemize}

\section{Results (pipeline verification)}

\subsection{Ranking}

\begin{center}
\begin{tabular}{@{}lcc@{}}
\toprule
Predictor & Spearman $\rho$ & Kendall $\tau_b$ \\
\midrule
predicted MACs          & 0.880 & 0.692 \\
predicted latency       & 0.765 & 0.580 \\
score, no resource term & 0.747 & 0.572 \\
\textbf{score, full}    & \textbf{0.746} & \textbf{0.556} \\
predicted memory        & 0.551 & 0.397 \\
capacity alone          & 0.512 & 0.402 \\
score, no capacity term & $-$0.684 & $-$0.501 \\
\bottomrule
\end{tabular}
\end{center}

The sign of the ablation matters more than the magnitudes. Removing the capacity
term makes the score \emph{anti-correlated} with accuracy: what remains only
prefers cheap configurations, which are the least accurate ones. The capacity
term therefore carries the accuracy signal and the resource term carries the
feasibility signal, which is precisely the claim the ablation is designed to
test.

\subsection{Selection under a fixed budget}

\begin{center}
\begin{tabular}{@{}lcccc@{}}
\toprule
Algorithm & Budget & Evaluated & Best accuracy & Regret \\
\midrule
random     & 4  & 4  & 0.456 & 0.094 \\
greedy     & 4  & 4  & 0.431 & 0.120 \\
heuristic  & 4  & 4  & 0.550 & 0.000 \\
random     & 12 & 12 & 0.521 & 0.030 \\
greedy     & 12 & 12 & 0.535 & 0.015 \\
heuristic  & 12 & 12 & 0.550 & 0.000 \\
exhaustive & all & 54 & 0.550 & 0.000 \\
\bottomrule
\end{tabular}
\end{center}

With a 400\,ms latency cap the picture changes, because feasibility starts to
bind: at budget 8 the heuristic reaches 0.422 with four feasible configurations
found, against three for random and two for greedy.

\subsection{Ranking under a hard constraint}

Once a constraint binds, the population that matters is the feasible one. At a
400\,ms cap (22 of 54 configurations feasible):

\begin{center}
\begin{tabular}{@{}lcc@{}}
\toprule
Predictor & $\rho$ (all 54) & $\rho$ (feasible only) \\
\midrule
predicted MACs          & $0.880$ & $0.853$ \\
predicted latency       & $0.765$ & $0.590$ \\
score, no resource term & $0.747$ & $0.641$ \\
\textbf{score, full}    & $-0.285$ & $\mathbf{0.656}$ \\
score, no capacity term & $-0.731$ & $-0.373$ \\
\bottomrule
\end{tabular}
\end{center}

On the full pool the score correlates \emph{negatively} with raw accuracy by
design: it ranks infeasible configurations last, and accuracy does not care about
the cap. Restricted to deployable configurations it recovers $\rho=0.66$ and
edges out the resource-free variant ($0.641$). Clamping each violation penalty at
one unit moved the full-pool figure from $-0.42$ to $-0.29$; the clamp is a
documented modelling choice, not a tuned parameter.

\subsection{Pareto front}

Twelve configurations are non-dominated in accuracy versus latency. The front
shows the expected staircase: small resolutions give cheap but weak models, and
the marginal accuracy per additional millisecond falls sharply above roughly
600\,ms.

\section{Limitations and threats to validity}

\begin{itemize}
  \item \textbf{Placeholder data.} The tables above come from a synthetic table.
        The synthetic truth is a smooth function of the same factors the score
        uses, so the heuristic looks strong here \emph{by construction}; only
        the hardware table can show genuine rank disagreement.
  \item \textbf{Simplified architecture model.} The analytic spec reproduces
        about 80--84\% of the published parameter counts and 93--94\% of the
        published GFLOPs. The compute match validates the part that drives the
        latency prediction; the parameter deficit is uniform, so it does not
        reorder the variants.
  \item \textbf{Rank correlation under a constraint.} Away from any cap the
        score ranks the pool well, but once a hard cap binds its correlation with
        raw accuracy over the \emph{full} pool is negative by construction, since
        it deliberately demotes infeasible configurations. Correlation must be
        reported on the feasible subset, alongside regret; the full-pool figure
        is not a meaningful statistic there.
  \item \textbf{Single measurement per configuration.} There are no seed-level
        repeats, so training-seed uncertainty cannot be estimated; random
        selection variability is reported separately.
  \item \textbf{Latency is not resource usage.} The reference FLOP estimate is
        not measured latency, and peak RSS is not per-inference memory.
\end{itemize}

\section{Reproducibility statement}

The repository is designed so that a fresh machine can regenerate every table
and figure:

\begin{verbatim}
pip install -r requirements.txt
python scripts/reproduce.py
\end{verbatim}

This runs the self-tests, the search comparison and the figure generation, and
writes \texttt{results/reproduce\_summary.txt}. The archived benchmark table is
committed, so no detector has to be run. Regenerating the measurements on real
hardware is a separate, explicitly optional step
(\texttt{scripts/benchmark\_pi.py}) that requires a Raspberry Pi and the optional
dependencies in \texttt{requirements-benchmark.txt}.

\begin{thebibliography}{9}
\setlength{\itemsep}{0pt}
\small
\bibitem{nsc} Zhu, Zhao and Lu. \emph{Neural Spectral Capacity: Measuring and
Designing Architectures from Network Specification Alone}.
\bibitem{lts} Javaheripi et al. \emph{LiteTransformerSearch: Training-free NAS for
Efficient Language Models}. NeurIPS 2022.
\bibitem{yolo} Ultralytics. \emph{YOLOv5}.
\bibitem{coco} Lin et al. \emph{Microsoft COCO}. ECCV 2014.
\bibitem{mp} Marchenko and Pastur. \emph{Distribution of eigenvalues for some sets
of random matrices}. 1967.
\bibitem{ort} ONNX Runtime, CPU execution provider.
\end{thebibliography}

\end{document}

